\documentclass[12pt,a4paper]{article}
\usepackage[utf8]{inputenc}
\usepackage[brazil]{babel}
\usepackage{graphicx}
\usepackage{amssymb, amsfonts, amsmath}
\usepackage{float}
\usepackage{enumerate}
\usepackage[top=2.5cm, bottom=2.5cm, left=1.25cm, right=1.25cm]{geometry}

\begin{document}
\pagestyle{empty}

\begin{center}
  \begin{tabular}{ccc}
    \begin{tabular}{c}
      \includegraphics[scale=0.25]{../../biblioteca/imagem/brasao-de-armas-brasil} \\
    \end{tabular} & 
    \begin{tabular}{c}
      Ministério da Educação \\
      Universidade Federal dos Vales do Jequitinhonha e Mucuri \\
      Faculdade de Ciências Sociais, Aplicadas e Exatas - FACSAE \\
      Departamento de Ciências Exatas - DCEX \\
      Disciplina: Introdução à Ciência da Computação \quad Semestre: 2023/2\\
      Prof. Dr. Luiz C. M. de Aquino\\
      Aluno(a):\rule{6cm}{0.1mm} \quad Data: \rule{0.5cm}{0.1mm}/\rule{0.5cm}{0.1mm}/\rule{1cm}{0.1mm}\\
    \end{tabular} &
    \begin{tabular}{c}
      \includegraphics[scale=0.25]{../../biblioteca/imagem/logo-ufvjm} \\
    \end{tabular}
  \end{tabular}
\end{center}

\begin{center}
 \textbf{Avaliação II}
\end{center}

\textbf{Instruções}
\begin{itemize}
 \item Todas as justificativas necessárias na solução de cada questão devem 
 estar presentes nesta avaliação;
 \item As respostas finais de cada questão devem estar escritas de caneta;
 \item Esta avaliação tem um total de 35,0 pontos.
\end{itemize}

\begin{enumerate}
  \item \textbf{[7,0 pontos]} Faça um programa que leia um vetor com 10 números reais.
  Em seguida, calcule a soma e o produto desses valores.
  
  \item \textbf{[7,0 pontos]} Faça um programa que leia 10 números reais e armazene em um vetor.
  Em seguida, calcule o desvio padrão nesse conjunto de números. Observação:
  considerando que o vetor seja \texttt{v}, o desvio padrão $\sigma$ é 
  calculado pela expressão:
  $$\sigma = \sqrt{\frac{1}{10} \sum_{i = 0}^9 (\texttt{v}[i] - m)^2}\textrm{,}$$
  onde $m$ é a média dos valores do vetor.
  
  \item \textbf{[7,0 pontos]} Faça um programa que leia um número natural $n$ e crie um vetor $\texttt{v}$
  com os algarismos desse número. Por exemplo, se $n = 1984$, então
  $\texttt{v} = [1,\,9,\,8,\,4]$.
  
  \item \textbf{[7,0 pontos]} Faça um programa que recebe um vetor $\texttt{v}$ com 10 números 
  inteiros e procure nele um certo número inteiro $n$ dado. Caso $n$ seja encontrado no vetor, 
  o programa deve imprimir a posição que ele está no vetor. Caso contrário, o programa deve 
  imprimir a mensagem que $n$ não foi encontrado.

  \item \textbf{[7,0 pontos]} Faça um programa que recebe um vetor $\texttt{v}$ com 10
  números inteiros e coloque em ordem crescente.

\end{enumerate}

\end{document}